
\subsection{例}
\label{subsec:3-7-6}

本小节给出一些简单而有启发性的弗洛尔上同调计算例子，并从几何上说明
弗洛尔上同调的存在性同第 3.6.3、3.7.3 小节讨论的势函数
$\mathfrak{PO}$ 之间的关系。此外，还将展示系数环
$\Lambda_{0,\mathrm{nov}}$ 与 $\Lambda_{\mathrm{nov}}$ 的选择，如何影响
弗洛尔上同调在哈密顿同痕下的不变性或依赖性。更多例子见
\textbf{[FOOO08I,08II,09I,09II]}。

考虑 $S^2\cong\mathbb CP^1$ 中一条嵌入圆周 $L$。以 $L$ 为边界有两个
简单全纯圆盘，记作 $D_u,D_\ell$。（它们实际上还有多重覆盖，但由次数理由
可知，这些多重覆盖对弗洛尔上同调没有贡献。）若两个圆盘的面积相同，例如
$L$ 是赤道时，则二者在
$\Pi(L)=\Pi(S^2;L)\cong\mathbb Z$ 中属于同一个类 $\beta$；记号见
定义 2.4.17。否则，它们在
$\Pi(S^2;L)\cong\mathbb Z\oplus\mathbb Z$ 中属于两个不同的类，分别记作
$\beta_u,\beta_\ell$。容易看出
\begin{equation}
 \mu_L(\beta)=\mu_L(\beta_u)=\mu_L(\beta_\ell)=2.
\tag{3.7.79}\label{eq:3-7-79}
\end{equation}
\liftIndex{I131}

以 $S^1_{\mathrm{eq}}$ 表示满足
$\Pi(S^1_{\mathrm{eq}})=\mathbb Z$ 的圆周，以
$S^1_{\mathrm{neq}}$ 表示满足
$\Pi(S^1_{\mathrm{neq}})=\mathbb Z\oplus\mathbb Z$ 的圆周。由
\eqref{eq:3-7-79} 可知，$S^2$ 中任意拉格朗日圆周 $L$ 都弱无障碍；
障碍类只出现在 $H^0(L;\mathbb Q)$ 中。

\par\medskip
\noindent\textbf{情形 I.}\enspace
\textit{$L^{(1)}$ 是 $L^{(0)}$ 的哈密顿形变。}
取哈密顿同痕 $\{\psi_t\}_{0\leq t\leq1}$，使得
$L^{(1)}=\psi_1(L^{(0)})$。

\par\medskip
\noindent\textbf{情形 I-a.}\enspace
\textit{$L^{(0)}=S^1_{\mathrm{eq}}$。}

第一，计算首个障碍类 $[o_1(S^1_{\mathrm{eq}})]$。如上所述，存在两个面积
相同的伪全纯圆盘 $D_u,D_\ell$。在 $S^2$ 上有一个反辛对合 $\tau$，它固定
$S^1_{\mathrm{eq}}$，并且 $\tau(D_u)=D_\ell$。定义障碍类时，要同时考虑
\[
 \mathcal M_1(S^1_{\mathrm{eq}};\beta),
 \qquad
 \mathcal M_1(S^1_{\mathrm{eq}};\tau_*\beta),
\]
即各带一个标记点的伪全纯圆盘模空间。因此，为了精确计算障碍类，需要分析
$\tau$ 如何作用于带标记点伪全纯圆盘模空间的定向；这在
\textbf{[FOOO09I]} 中作了研究。由那里可知，
\[
 \tau_*\colon
 \mathcal M_1(S^1_{\mathrm{eq}};\beta)
 \longrightarrow
 \mathcal M_1(S^1_{\mathrm{eq}};\tau_*\beta)
\]
保持定向，因为 $\mu_L(\beta)=2$，且这里只有一个标记点。所以
$D_u,D_\ell$ 对障碍类
$[o_1(S^1_{\mathrm{eq}})]=[o(S^1_{\mathrm{eq}};\bar\beta)]$
的贡献相同。这里
\[
 \bar\beta=[\beta]=[\tau_*\beta]\in\Pi(S^1_{\mathrm{eq}}).
\]
故
\begin{equation}
 [o(S^1_{\mathrm{eq}};\bar\beta)]
 =2[S^1_{\mathrm{eq}}]\in H^0(S^1;\mathbb Q).
\tag{3.7.80}\label{eq:3-7-80}
\end{equation}

第二，不存在马斯洛夫指标为负的伪全纯圆盘；事实上，这正是单调情形。因此
弗洛尔上同调可以定义，即
$S^1_{\mathrm{eq}}\subset S^2$ 弱无障碍。而且
$0\in\widehat{\mathcal M}_{\mathrm{weak}}(S^1_{\mathrm{eq}})$。本小节余下
部分均从记号中略去 $0$。

第三，用博特--莫尔斯理论计算
$HF(L^{(0)},L^{(0)};\Lambda_{0,\mathrm{nov}})$
\translatornote{原书在这一引导句中写
$HF(L^{(0)},L^{(0)};\Lambda_{\mathrm{nov}})$；但紧接的
\eqref{eq:3-7-81} 以及随后对两个系数环的对照都以
$\Lambda_{0,\mathrm{nov}}$ 为这里的计算对象，译文据此恢复。}。这里
$L^{(0)}=S^1_{\mathrm{eq}}$。当 $\beta\ne0$ 时，上边界映射
$\mathfrak m_{1,\beta}$ 由模空间
$\mathcal M_2(L^{(0)};\beta;J)$ 诱导。依 \textbf{[FOOO09I]} 的结果，
\[
 ev_*\bigl([\mathcal M_2(L^{(0)};\beta;J)]\bigr)
 =
 -ev_*\bigl([\mathcal M_2(L^{(0)};\tau_*\beta;J)]\bigr).
\]
这是因为此时有两个标记点，所以 $\tau_*$ 反转定向。因而
\begin{equation}
 HF(L^{(0)},L^{(0)};\Lambda_{0,\mathrm{nov}})
 \cong\Lambda_{0,\mathrm{nov}}\oplus\Lambda_{0,\mathrm{nov}}.
\tag{3.7.81}\label{eq:3-7-81}
\end{equation}
换一种说法，第 6 章将构造的谱序列在 $E_2$ 页退化。

下面把 \eqref{eq:3-7-81} 同
$HF(L^{(1)},L^{(0)};\Lambda_{\mathrm{nov}})$ 作比较，其中
$L^{(1)}=\psi_1(L^{(0)})$ 与 $L^{(0)}$ 哈密顿同痕。取哈密顿同痕使
\[
 \psi_1(L^{(0)})\cap L^{(0)}=\{p,q\}.
\]
一般而言，应注意
$\#(\psi_1(S^1_{\mathrm{eq}})\cap S^1_{\mathrm{eq}})\geq2$。从 $p$ 到
$q$ 以及从 $q$ 到 $p$ 的极小面积连接轨道模空间各有两个分支。

另一方面，由弗洛尔上同调的哈密顿同痕不变性及
\eqref{eq:3-7-81}，有
\begin{equation}
 HF(\psi_1(S^1_{\mathrm{eq}}),S^1_{\mathrm{eq}};
 \Lambda_{\mathrm{nov}})
 \cong\Lambda_{\mathrm{nov}}\oplus\Lambda_{\mathrm{nov}}.
\tag{3.7.82.1}\label{eq:3-7-82-1}
\end{equation}
因此，上述两个分支的贡献彼此抵消。（也可以把它同
$L^{(0)}=S^1$ 上一个恰有两个临界点 $\{p,q\}$ 的莫尔斯函数之梯度轨线
模空间比较，从而直接证明这一点。）非极小能量伪全纯条带的贡献由次数理由
为零。于是
\begin{equation}
 HF(\psi_1(S^1_{\mathrm{eq}}),S^1_{\mathrm{eq}};
 \Lambda_{0,\mathrm{nov}})
 \cong\Lambda_{0,\mathrm{nov}}\oplus\Lambda_{0,\mathrm{nov}}.
\tag{3.7.82.2}\label{eq:3-7-82-2}
\end{equation}

\par\medskip
\noindent\textbf{情形 I-b.}\enspace
\textit{$L^{(0)}=S^1_{\mathrm{neq}}$。}

这时，以 $L^{(0)}=S^1_{\mathrm{neq}}$ 为边界的两个伪全纯圆盘
$D_u,D_\ell$ 具有不同的辛面积。与生成
$\Pi(S^2;S^1_{\mathrm{neq}})\cong\mathbb Z\oplus\mathbb Z$ 的
$\beta_u,\beta_\ell$ 对应，有障碍类
$[o(S^1_{\mathrm{neq}};\beta_u)]$ 与
$[o(S^1_{\mathrm{neq}};\beta_\ell)]$。可以看出
\begin{equation}
 [o(S^1_{\mathrm{neq}};\beta_u)]
 =
 [o(S^1_{\mathrm{neq}};\beta_\ell)]
 =
 [S^1_{\mathrm{neq}}]\in H^0(S^1;\mathbb Q).
\tag{3.7.83}\label{eq:3-7-83}
\end{equation}

来计算弗洛尔上同调
$HF(\psi_1(S^1_{\mathrm{neq}}),S^1_{\mathrm{neq}})$。这里要强调：一般
而言，取 $\Lambda_{0,\mathrm{nov}}$ 系数的弗洛尔上同调依赖于哈密顿微分
同胚 $\psi_1$，而取 $\Lambda_{\mathrm{nov}}$ 系数时则总与该哈密顿微分
同胚无关。这一点将在第 5.3 节作为定理 5.3.14 证明；这里用本例具体说明
这一现象。

首先，存在幅度充分大的哈密顿同痕
$\{\psi_t\}_{0\leq t\leq1}$，使
\[
 \psi_1(S^1_{\mathrm{neq}})\cap S^1_{\mathrm{neq}}=\varnothing.
\]
在这种情形下，
\[
 HF(\psi_1(S^1_{\mathrm{neq}}),S^1_{\mathrm{neq}};
 \Lambda_{0,\mathrm{nov}})=0.
\]
取 $\Lambda_{\mathrm{nov}}$ 系数时显然也为零。

\begin{figure}[htbp]
  \centering
  \setlength{\unitlength}{0.78mm}
  \begin{picture}(118,98)
    \put(40,12){\line(1,0){44}}
    \put(34,18){\line(1,0){56}}
    \put(28,24){\line(1,0){68}}
    \put(24,30){\line(1,0){76}}
    \put(22,36){\line(1,0){80}}
    \qbezier(39,75)(50,60)(61,45)
    \qbezier(48,82)(61,64)(74,46)
    \qbezier(59,87)(73,68)(87,49)
    \qbezier(71,86)(85,70)(99,54)
    \qbezier(24,55)(31,49)(38,43)
    \qbezier(27,62)(37,52)(47,42)
    \qbezier(31,68)(44,55)(57,42)
    \qbezier(78,41.8)(78.9,40.6)(79.8,39.4)
    \qbezier(86,42.2)(86.7,41.25)(87.4,40.3)
    \qbezier(98.3,45.9)(99.55,44.8)(100.8,43.7)
    \qbezier(100.2,48.1)(101.95,45.75)(103.7,43.4)
    \put(62,48){\circle{84}}
    \qbezier(20,47)(57,38)(104,44)
    \qbezier(29,70)(45,49)(70,42)
    \qbezier(70,42)(91,33)(104,53)
    \qbezier[34](29,70)(55,67)(78,58)
    \qbezier[24](78,58)(94,56)(104,53)
    \qbezier[34](20,47)(49,61)(78,58)
    \qbezier[24](78,58)(94,52)(104,44)
    \put(70,42){\circle*{1.6}}
    \put(78,58){\circle*{1.6}}
    \put(69,36){$p$}
    \put(79,61){$q$}
    \put(61,72){$A$}
    \put(59,22){$B$}
    \put(16,53){$C$}
    \put(106,47){$D$}
    \put(1,43){$L^{(0)}$}
    \put(2,72){$\psi_1(L^{(0)})$}
  \end{picture}
  \par\smallskip
  \renewcommand{\thefigure}{3.7.10}
  \refstepcounter{figure}
  \textbf{\figurename~\thefigure}\label{fig:3-7-10}
\end{figure}

另一方面，考虑
$\psi_1(L^{(0)})\cap L^{(0)}=\{p,q\}$ 的情形。这时
\[
 \mathbb CP^1\setminus
 \bigl(\psi_1(L^{(0)})\cup L^{(0)}\bigr)
\]
有四个连通分支，分别以 $A,B,C,D$ 表示其辛面积。因为 $\psi_1$ 是哈密顿
微分同胚，所以 $C=D$ 或 $A=B$。（由于 $L^{(0)}$ 不是赤道，
$A=B$ 与 $C=D$ 同时成立的情形不在这里。）见图~\ref{fig:3-7-10}。

面积为 $A$ 的从 $p$ 到 $q$ 的连接轨道模空间，其定向同面积为 $B$ 的对应
模空间相反；类似地，面积为 $C$ 的从 $q$ 到 $p$ 的模空间，其定向同面积
为 $D$ 的对应模空间相反。这可由此处模空间的定向同情形 I-a 中的定向一致
来证明；也可以约化到莫尔斯（诺维科夫）理论来证明。当 $C=D$ 时，从 $p$
到 $q$ 的两条连接轨道在覆盖变换群
$\Pi(S^2;S^1_{\mathrm{neq}})$ 的作用下落在不同轨道中，而从 $q$ 到 $p$
的两条连接轨道彼此抵消。$A=B$ 的情形类似。因此，上边界算子 $\delta$
由下式给出：
\begin{equation}
 [p]\longmapsto T^A[q]-T^B[q],
 \qquad
 [q]\longmapsto0,
\tag{3.7.84.1}\label{eq:3-7-84-1}
\end{equation}
或者
\begin{equation}
 [p]\longmapsto0,
 \qquad
 [q]\longmapsto T^C[p]e-T^D[p]e.
\tag{3.7.84.2}\label{eq:3-7-84-2}
\end{equation}
这里把 $[q]$ 的次数取成比 $[p]$ 大 $1$，并回忆依本书约定，不定元 $e$ 的
次数为 $2$。伪全纯圆盘模空间其他分支对上边界算子的贡献由次数理由为零。
于是 $\delta\circ\delta=0$，故
\[
 HF(\psi_1(S^1_{\mathrm{neq}}),S^1_{\mathrm{neq}};
 \Lambda_{0,\mathrm{nov}})
\]
也可以定义，而且在当前情形下不为零。它是一个挠
$\Lambda_{0,\mathrm{nov}}$-模；事实上，若 $C<D$，它就是
$\Lambda_{0,\mathrm{nov}}/T^C\Lambda_{0,\mathrm{nov}}$。这个例子表明，
$HF(\psi_1(S^1_{\mathrm{neq}}),S^1_{\mathrm{neq}};
\Lambda_{0,\mathrm{nov}})$ 的确依赖于哈密顿微分同胚。不过，取
$\Lambda_{\mathrm{nov}}$ 系数时，有
\begin{equation}
 HF(\psi_1(S^1_{\mathrm{neq}}),S^1_{\mathrm{neq}};
 \Lambda_{\mathrm{nov}})=0.
\tag{3.7.85}\label{eq:3-7-85}
\end{equation}
这是因为，当 $C=D$ 时，$T^A-T^B$ 在 $\Lambda_{\mathrm{nov}}$ 中可逆；
当 $A=B$ 时，$T^C-T^D$ 在 $\Lambda_{\mathrm{nov}}$ 中可逆。事实上，若
$A<B$，则
\[
 (T^A-T^B)^{-1}
 =T^{-A}(1-T^{B-A})^{-1}
 =T^{-A}\sum_{j=0}^{\infty}T^{(B-A)j}.
\]
这一结果同
$\psi_1(S^1_{\mathrm{neq}})\cap S^1_{\mathrm{neq}}=\varnothing$
时取 $\Lambda_{\mathrm{nov}}$ 系数的弗洛尔上同调一致。

\par\medskip
\noindent\textbf{情形 II.}\enspace
\textit{$L^{(1)}$ 不是 $L^{(0)}$ 的哈密顿形变。}

考虑 $L^{(1)}\cap L^{(0)}=\{p,q\}$ 的情形。同情形 I-b 一样，可以看出
上边界算子应当是
\[
 [p]\longmapsto T^A[q]-T^B[q],
 \qquad
 [q]\longmapsto T^C[p]e-T^D[p]e.
\]
但由于 $L^{(1)}$ 不是 $L^{(0)}$ 的哈密顿形变，此时
$A\ne B$ 且 $C\ne D$。所以 $\delta\circ\delta\ne0$，弗洛尔上同调也就
没有定义。

\begin{remark}{3.7.86}\label{rem:3-7-86}
在第 3.6.1 小节的无障碍情形和第 3.6.3 小节的弱无障碍情形中，我们定义了
弗洛尔上同调 $HF(L,b;\Lambda_{0,\mathrm{nov}})$，见定义 3.6.11 和
定义 3.6.33 (1)。另一方面，定义 3.7.74 在第 3.7.5 小节把
\[
 HF((L,b),(L,b);\Lambda_{0,\mathrm{nov}})
\]
定义为博特--莫尔斯情形下拉格朗日子流形对子之弗洛尔上同调的一个特例。
命题 3.7.75 说明，这两个上同调在 $\Lambda_{0,\mathrm{nov}}$ 上同构。
然而，正如上例所示，对哈密顿微分同胚 $\psi\ne\mathrm{id}$，它们一般并不
在 $\Lambda_{0,\mathrm{nov}}$ 上同构于对子
$((L,b),(\psi(L),\psi_*(b)))$ 的弗洛尔上同调；在
$\Lambda_{\mathrm{nov}}$ 上则同构。这将在第 4.1 节作为定理 4.1.5 证明
\translatornote{原书把定理 4.1.5 的位置写成“第 5 章第 5.3 节”；
定理编号及该定理的实际位置均指向第 4.1 节，译文据此修正。}。此外，
\[
 HF((\psi(L),\psi_*(b)),(L,b);\Lambda_{0,\mathrm{nov}})
\]
对 $\psi$ 的依赖体现在弗洛尔上同调的挠部分中。第 6.5.3 小节将研究这个
挠部分，并把它用于研究某些哈密顿同痕的霍费尔距离。
\end{remark}

\subsection{弗洛尔上同调的乘法结构}
\label{subsec:3-7-7}

本小节证明若干纯代数结果；第 3.8.10 小节将用它们证明定理 B 的 (B.2)
和定理 G 的 (G.5)。\liftIndex{I130}

\begin{proposition}{3.7.87}\label{prop:3-7-87}
\begin{enumerate}[label=\textup{(\arabic*)},leftmargin=2.8em]
 \item 设 $C$ 是滤过 $A_\infty$ 代数，且
 $b_i\in\widehat{\mathcal M}(C)$（$i=0,1,2$）。则存在运算
 \begin{equation}
  \mathfrak m_2\colon
  H(C,\delta_{b_0,b_1})\otimes H(C,\delta_{b_1,b_2})
  \longrightarrow H(C,\delta_{b_0,b_2}),
 \tag{3.7.88}\label{eq:3-7-88}
 \end{equation}
 它在相差符号的意义下结合；也就是说
 \translatornote{原书在 \eqref{eq:3-7-89} 中把指数写成
 $\deg' x$。由微分分次代数特例 \eqref{eq:3-2-5} 直接核算，所需指数是
 $\deg x=\deg' x+1$；也只有这一符号才使紧接的
 \eqref{eq:3-7-92} 成为结合的杯积。译文据此恢复。}，
 \begin{equation}
  \mathfrak m_2(\mathfrak m_2(x,y),z)
  =
  (-1)^{\deg' x+1}\mathfrak m_2(x,\mathfrak m_2(y,z)).
 \tag{3.7.89}\label{eq:3-7-89}
 \end{equation}

 \item 若 $C$ 含幺，则对
 $b_i\in\widehat{\mathcal M}_{\mathrm{weak}}(C)$（$i=0,1,2$），只要
 \[
  \mathfrak{PO}(b_0)=\mathfrak{PO}(b_1)=\mathfrak{PO}(b_2),
 \]
 同一结论仍成立。
 \translatornote{原书第 (2) 款只写
 $\mathfrak{PO}(b_1)=\mathfrak{PO}(b_0)$，漏掉 $b_2$。但
 \eqref{eq:3-7-88} 中涉及的三个上同调都必须定义；由命题 3.7.17，
 这要求三个势函数值相等。}
\end{enumerate}
\end{proposition}

\begin{proof}
令 $x_{ij}\in C$，并置
\begin{equation}
 \widetilde{\mathfrak m}_2(x_{ij},x_{jk})
 =
 \mathfrak m(e^{b_i}x_{ij}e^{b_j}x_{jk}e^{b_k}).
\tag{3.7.90}\label{eq:3-7-90}
\end{equation}
由 $A_\infty$ 公式容易看出
\begin{equation}
\begin{aligned}
 -(\delta_{b_i,b_k}\circ\widetilde{\mathfrak m}_2)
 (x_{ij},x_{jk})
 ={}&
 \widetilde{\mathfrak m}_2(\delta_{b_i,b_j}x_{ij},x_{jk})\\
 &+(-1)^{\deg' x_{ij}}
 \widetilde{\mathfrak m}_2(x_{ij},\delta_{b_j,b_k}x_{jk}).
\end{aligned}
\tag{3.7.91}\label{eq:3-7-91}
\end{equation}
因此，\eqref{eq:3-7-90} 诱导 \eqref{eq:3-7-88}。再置
\[
 \widetilde{\mathfrak m}_3(x_{hi},x_{ij},x_{jk})
 =
 \mathfrak m(
 e^{b_h}x_{hi}e^{b_i}x_{ij}e^{b_j}x_{jk}e^{b_k}).
\]
若 $\delta_{b_u,b_v}(x_{uv})=0$，则有
\translatornote{原书在 \eqref{eq:3-7-91} 左端把第二个输入误写成
$x_{ik}$；又在后一条同伦公式左端把第三个输入写成 $x_{ik}$，并把第二项
的符号指数写成 $\deg' x_{ij}$。定义 \eqref{eq:3-7-90} 及各项的首尾指标
分别强制两个输入为 $(x_{ij},x_{jk})$、三个输入为
$(x_{hi},x_{ij},x_{jk})$；余导子越过第一个输入时的符号则是
$(-1)^{\deg' x_{hi}}$。译文据此恢复。}
\[
\begin{aligned}
 -(\delta_{b_h,b_k}\circ\widetilde{\mathfrak m}_3)
 (x_{hi},x_{ij},x_{jk})
 ={}&
 \widetilde{\mathfrak m}_2(
   \widetilde{\mathfrak m}_2(x_{hi},x_{ij}),x_{jk})\\
 &+(-1)^{\deg' x_{hi}}
 \widetilde{\mathfrak m}_2(
   x_{hi},\widetilde{\mathfrak m}_2(x_{ij},x_{jk})).
\end{aligned}
\]
由此得到 \eqref{eq:3-7-89}。弱版本的证明相同。
\end{proof}

注意，\eqref{eq:3-7-89} 在 $H(C,\delta_{b,b})$ 上诱导如下乘积：
\begin{equation}
 x\cup y
 =
 (-1)^{\deg x(\deg y+1)}\mathfrak m_2(x,y).
\tag{3.7.92}\label{eq:3-7-92}
\end{equation}

\begin{proposition}{3.7.93}\label{prop:3-7-93}
\begin{enumerate}[label=\textup{(\arabic*)},leftmargin=2.8em]
 \item 设 $C_1,C_0$ 是滤过 $A_\infty$ 代数，$D$ 是
 $(C_1,C_0)$ 滤过 $A_\infty$ 双模，并且
 $b_i\in\widehat{\mathcal M}(C_i)$（$i=0,1$）。则
 $H(D,\delta_{b_1,b_0})$ 是左
 $H(C_1,\delta_{b_1,b_1})$、右
 $H(C_0,\delta_{b_0,b_0})$ 的双模
 \translatornote{原书在第 (1) 款的左右作用代数以及证明中都写
 $H(C,\delta_{b_i,b_i})$，但命题只给出了 $C_1,C_0$，并未定义 $C$。
 依双模的左右代数及 $b_i$ 的所属关系，译文分别恢复为 $C_1,C_0$。}。

 \item 若 $C_i$ 与 $D$ 都含幺，则对满足
 \[
  b_i\in\widehat{\mathcal M}_{\mathrm{weak}}(C_i),
  \qquad
  \mathfrak{PO}(b_1)=\mathfrak{PO}(b_0),
 \]
 的 $b_i$，同一结论仍成立。
\end{enumerate}
\end{proposition}

\begin{proof}
定义
\[
 \widetilde{\mathfrak n}_{1,0}\colon C_1\otimes D\longrightarrow D,
 \qquad
 \widetilde{\mathfrak n}_{1,0}(x,v)
 =
 \mathfrak n(e^{b_1}xe^{b_1},v,e^{b_0}).
\]
容易看出，$\widetilde{\mathfrak n}_{1,0}$ 在
$H(D,\delta_{b_1,b_0})$ 上诱导左
$H(C_1,\delta_{b_1,b_1})$-模结构。另一侧的证明类似。
\end{proof}
